\scanpage{7}
\begin{center}Sur les groupes de Galois attachés aux groupes $p$-divisibles\hfill121\end{center}

\textit{vectoriel de dimension finie, et soit $H$ un sous-groupe algébrique réductif de $\mathbf{GL}_V$. Pour qu'un point de $\mathbf{GL}_V$, à valeurs dans une extension de $E$, soit un point de $H$, il faut et il suffit que, pour tout couple $(r,s)$, il laisse invariants les éléments de $T_{r,s}(V)$ invariants par $H$.}

(En d'autres termes, un groupe réductif est déterminé par ses invariants tensoriels.)

Comme c'est là un résultat bien connu, je me bornerai à en esquisser la démonstration. On se ramène tout de suite au cas où le corps de base est algébriquement clos, ce qui permet d'identifier le groupe à l'ensemble de ses points rationnels; on peut aussi supposer que $H$ est contenu dans le groupe unimodulaire $\mathbf{SL}_V$ (en effet, si $n=\dim.V$, on remplace $V$ par $V\oplus\bigwedge^n V^*$). Soit alors $g$ un élément de $\mathbf{GL}_V$, vérifiant la condition du lemme, et n'appartenant pas à $H$. Du fait que $H$ est contenu dans $\mathbf{SL}_V$, $H$ et $Hg$ sont fermés dans $\operatorname{End}(V)$; puisque $H$ est réductif, il existe donc une fonction polynôme $f$ sur $\operatorname{End}(V)$, invariante par multiplication à gauche par $H$, et prenant les valeurs 0 sur $H$ et 1 sur $Hg$ (cf. [4], p. 29, cor. 1.2). Mais la représentation de $H$ dans $\operatorname{End}(V)$ par multiplication à gauche est isomorphe à la somme directe de $n$ copies de $V$. Comme $f$ appartient à l'algèbre symétrique du dual de $\operatorname{End}(V)$, l'hypothèse faite sur $g$ montre que $f$ est invariante par $g$. Mais c'est absurde, puisque $f(1)=0$ et $f(g)=1$.

\begin{center}\textbf{\S\ 2. Modules galoisiens du type de Hodge-Tate}\end{center}

A partir de maintenant, $K$ désigne un corps muni d'une valuation discrète, et complet pour la topologie définie par cette valuation. On note $A$ l'anneau des entiers de $K$, et $k=A/\mathfrak m$ le corps résiduel correspondant. On suppose $K$ de caractéristique zéro, et $k$ parfait de caractéristique $p$. On note $\overline K$ une clôture algébrique de $K$, et $C$ sa complétion. Le groupe de Galois $\operatorname{Gal}_K$ de $\overline K/K$ opère sur $C$ (cf. [9], \S\ 3).

Soit $\mathbf U_p=\mathbf Z_p^*$ le groupe des unités de $\mathbf Q_p$, et soit
\[
\chi:\operatorname{Gal}_K\to\mathbf U_p
\]
le caractère de $\operatorname{Gal}_K$ donnant l'action de ce groupe sur les racines $p^n$-èmes de l'unité. Par définition, on a
\[
sz=z^{\chi(s)}
\]
pour tout $s\in\operatorname{Gal}_K$ et pour toute racine de l'unité $z$ d'ordre une puissance de $p$.

Soit maintenant $X$ un $C$-espace vectoriel de dimension finie muni d'une

9\qquad Local Fields
\margin{6}
